\chapter{Introducere}
\label{cap:introducere}

\section{Tipul lucrării și încadrarea temei}

Lucrarea de față este o lucrare de licență aplicativă, încadrată în zona aplicațiilor web full-stack
și a sistemelor de gestiune a datelor pentru domeniul fitness. Tema nu constă doar într-un jurnal de
antrenamente, ci într-o platformă completă în care utilizatorul își poate planifica antrenamentele,
le poate executa, le poate analiza după ce au fost realizate și poate descoperi programe create de
alți utilizatori.

Din punct de vedere tehnic, proiectul combină mai multe subprobleme care apar într-un sistem real:
autentificare sigură, autorizare pe resurse, date istorice care nu trebuie alterate, operații
concurente, căutare full-text, marketplace, analiză automată a structurii unui program și acces
controlat pentru antrenori. Tocmai această combinație este partea interesantă: fiecare funcționalitate
trebuie să fie utilă pentru utilizator, dar și corectă în raport cu restul sistemului.

\section{Prezentarea generală a temei}

Un utilizator al aplicației începe prin a-și construi biblioteca de lucru: exerciții oficiale și
exerciții proprii, rutine de încălzire sau încheiere, săli și echipamente. Pe baza acestora creează
programe de antrenament împărțite pe zile. O zi de program conține exerciții ordonate, valori
planificate (serii, repetări, greutate, pauză) și, opțional, rutine atașate.

În momentul în care utilizatorul pornește un antrenament, programul devine o sesiune live. Aici apar
datele istorice propriu-zise: seriile efectuate, greutatea, repetările, durata, RPE-ul, notele și
echipamentul folosit. După finalizare, sesiunea intră în istoric și este folosită în analizele de
progres. Sistemul calculează volum în timp, antrenamente pe săptămână, distribuția pe grupe musculare
și estimări de progres la forță.

Pe lângă fluxul individual, aplicația are un marketplace de programe publice. Un utilizator poate
publica un program, iar ceilalți îl pot vota, salva sau copia în biblioteca lor privată. Există și un
mod de coaching: un utilizator cu rol de antrenor poate vedea istoricul și analizele unui client numai
dacă există o relație activă acceptată de client. Acest acces este intenționat limitat la citire, ca
să nu amestece colaborarea cu dreptul de a modifica datele personale ale clientului.

\section{Problema urmărită}

Într-o aplicație de acest tip, problema cea mai importantă nu este doar să se poată salva antrenamente,
ci să se poată avea încredere în ce s-a salvat. Datele de planificare sunt, prin natura lor,
editabile: un utilizator poate redenumi un exercițiu, poate modifica un program, poate șterge o sală
sau poate schimba echipamentul disponibil. În schimb, istoricul trebuie să descrie ce s-a întâmplat
la momentul respectiv. Dacă o sesiune veche s-ar schimba pentru că programul original a fost editat,
aplicația ar pierde sensul de jurnal.

De aici rezultă prima direcție a lucrării: un model de date cu instantanee pentru istoricul
antrenamentelor. La pornirea sesiunii, aplicația copiază datele relevante din planificare în tabele de
sesiune. Cheile către datele sursă rămân doar ca legături auxiliare și sunt anulabile, în timp ce
valorile afișate în istoric sunt cele copiate în momentul execuției.

A doua direcție este securitatea. Aplicația este multi-user și include resurse personale, relații
antrenor-client și endpoint-uri administrative. De aceea, fiecare operație trebuie să știe cine este
utilizatorul curent și ce are voie să vadă sau să modifice. Sistemul folosește JWT, token-uri de
reîmprospătare opace rotite, interogări scoping după proprietar și reguli de acces care întorc 404
pentru resursele altui utilizator, pentru a evita scurgerile de informație de tip IDOR.

A treia direcție este regăsirea informației. Pe măsură ce crește numărul de programe și sesiuni,
utilizatorul are nevoie de căutare după exerciții, grupe musculare, săli, echipamente, note și
filtre. În proiect, căutarea este implementată cu OpenSearch ca model de citire derivat din
PostgreSQL, dar lucrarea nu tratează OpenSearch ca soluție implicită. Ea include și o comparație
practică împotriva unei variante PostgreSQL bine indexate, pentru a vedea când merită complexitatea
unui motor separat.

\section{Contribuția proprie}

Contribuția proprie constă în proiectarea și implementarea unei platforme funcționale, nu doar a unui
prototip izolat. Principalele elemente realizate sunt:

\begin{itemize}
  \item un backend Spring Boot organizat pe funcționalități, cu separare între domeniu, servicii,
        repository-uri și controllere;
  \item o schemă PostgreSQL gestionată prin Flyway, cu constrângeri de integritate, indecși unici
        parțiali, ștergere logică și reguli de concurență la nivel de bază de date;
  \item un model de sesiuni cu instantanee, astfel încât istoricul rămâne corect după editarea sau
        ștergerea datelor de planificare;
  \item autentificare cu JWT, rotația token-urilor de reîmprospătare și protecție față de accesul
        neautorizat la resursele altor utilizatori;
  \item funcționalități complete pentru exerciții, rutine, săli, echipamente, programe, sesiuni live,
        istoric, analize, marketplace și coaching;
  \item un analizor determinist al structurii programelor, bazat pe reguli explicabile, fără a pretinde
        că oferă sfat medical;
  \item un model de citire OpenSearch pentru căutare full-text peste programe și istoric, cu filtre,
        fațete, evidențiere, toleranță la greșeli și verificare suplimentară în PostgreSQL;
  \item o evaluare experimentală OpenSearch vs. PostgreSQL pe un set de date generat până la două
        milioane de documente.
\end{itemize}

\section{Structura lucrării}

Capitolul~\ref{cap:preliminarii} introduce noțiunile folosite în proiect: aplicații web, modelarea
datelor istorice, autentificarea cu token-uri, controlul accesului, căutarea full-text și principiile
folosite de analizorul de programe. Capitolul~\ref{cap:contributie} descrie implementarea efectivă a
platformei, urmărind fluxul de utilizare și apoi deciziile tehnice: modelul de date, securitatea,
funcționalitățile, coaching-ul, căutarea, benchmark-ul și testarea. Capitolul~\ref{cap:concluzii}
rezumă rezultatele și discută direcțiile de dezvoltare rămase. Detaliile care ar încărca textul
principal, cum ar fi schema bazei de date, fragmente de cod, tabelele complete de benchmark și
contractul API, sunt incluse în anexe.
